\chapter{总结}\label{chap:summary}

本次课程设计围绕机器人避障寻径问题，完成了从基础算法实现、复杂场景设计，到多算法比较和动态可视化展示的一整套实验流程。

具体来说，本文首先实现了 BFS 和 A* 两种基础算法；随后自主设计了一个 30x30 的复杂障碍场景；在此基础上补充实现了改进 A*、Q-Learning、人工势场法和 Dijkstra；最后输出了静态结果图、训练曲线和 GIF 动态展示文件。

实验结果表明：
\begin{itemize}
  \item BFS 和 Dijkstra 能稳定得到最短路径，但搜索范围较大；
  \item A* 具有更高的搜索效率；
  \item 改进 A* 能在路径长度不变的前提下降低综合代价；
  \item Q-Learning 经过训练后能够得到可行路径；
  \item 人工势场法在复杂环境中容易受到局部极小值影响。
\end{itemize}

总体来看，本次课程设计让我对图搜索算法、强化学习方法以及路径规划中的可视化展示有了更直观的理解。
